


\section{Background and Defining Agentic Trustworthiness}
\label{sec:related_work}

\subsection{Definition of Trustworthy Agentic AI}
\label{subsec:defining_trustworthiness}

Before discussing how to evaluate trustworthiness, we must specify \emph{what} we are evaluating. Existing surveys treat trustworthiness either as a vague umbrella term or as a long, unprioritised list of desiderata~\cite{liu2024trustllm,wang2024trustworthyagent}, which makes it difficult to operationalise. Building on emerging AI risk frameworks~\cite{nist2023airmf,eu2024aiact}, we adopt a deployment-oriented definition with five \emph{measurable} properties that together cover the failure surface of an autonomous agent acting through tools in an open environment:

\begin{itemize}[leftmargin=*]
\item \textbf{(P1) Reliability.} Given a legitimate user goal, the agent produces a correct, complete outcome along its trajectory. Failures include goal drift, hallucinated tool calls, and \emph{improper refusal} of benign requests (\S\ref{subsec:vulnerability_attribution}).
\item \textbf{(P2) Robustness.} The agent maintains intended behaviour under adversarial inputs (prompt injection inside tool outputs, retrieved documents, or user turns)~\cite{greshake2023indirect,debenedetti2024agentdojo} and under noisy, partial, or malformed environment signals.
\item \textbf{(P3) Safety.} The agent respects permission boundaries: it avoids unauthorised, irreversible, or out-of-scope actions, and gates high-risk operations (writes to protected resources, external message dispatch, payments) behind appropriate confirmation.
\item \textbf{(P4) Social-Ethical Alignment.} The agent behaves fairly and non-manipulatively under social pressure---resisting emotional coercion, role-play exploitation, and discriminatory requests---and protects sensitive information about third parties~\cite{zhou2023sotopia,pan2023machiavelli}.
\item \textbf{(P5) Operational Integrity.} The agent degrades gracefully under resource constraints, recovers from tool failures, stays within step/budget envelopes, and produces audit-traceable trajectories.
\end{itemize}

We treat P1--P5 as the \emph{evaluation targets} of HAAF, not metrics layered on existing benchmarks. Three consequences follow: trustworthiness is \emph{multi-dimensional} (no single scalar suffices, since two systems with the same aggregate score may have incomparable profiles), \emph{distributional} (each property must be assessed over a representative scenario distribution), and \emph{deployment-conditioned} (it is a property of the full deployed system---weights + prompts + tools + guardrails---not the model in isolation~\cite{ganguli2022redteaming}). This definition governs the rest of the paper: \S\ref{subsec:overview}--\ref{subsec:layer4} map each layer to a subset of P1--P5 (Table~\ref{tab:layer_property}), the taxonomy in \S\ref{subsec:vulnerability_attribution} assigns failures to properties, and the profile in \S\ref{subsec:profile_comparison} reports P1--P5 separately.

\subsection{Benchmark Fragmentation and the Representativeness Gap}
\label{subsec:gap}

Current benchmark ecosystems have substantially advanced the evaluation of agentic systems by operationalising important local notions of competence and safety~\cite{mohammadi2025agenteval}. Task-centric benchmarks such as AgentBench, WebArena, and SWE-bench evaluate interactive task completion and software engineering~\cite{liu2023agentbench,zhou2023webarena,jimenez2023swebench}; HaluEval and JailbreakBench probe hallucination and adversarial robustness~\cite{li2023halueval,chao2024jailbreakbench}; API-Bank tests tool-augmented behaviour~\cite{li2023apibank}; and more recent suites---GAIA, WorkArena, OSWorld, $\tau$-bench, ToolSandbox, AgentDojo---extend coverage to general assistants, enterprise software, computer-use, and prompt-injection robustness~\cite{mialon2023gaia,drouin2024workarena,xie2024osworld,yao2024taubench,lu2024toolsandbox,debenedetti2024agentdojo}. These benchmarks provide valuable signals, but each is tied to a particular evaluation slice and a particular subset of P1--P5.

The problem is not that benchmarks are individually uninformative, but that they are collectively incomplete in three ways: \emph{capability-isolated} (coding, planning, hallucination, safety tested separately when real failures emerge from their interaction); \emph{environment-simplified} (clean tools and stable interfaces vs.\ real noise, latency, partial observability, cascading effects); and \emph{socially decontextualised} (delegation, manipulation, bias, risk asymmetry abstracted away). Combined, they reveal success under one pressure at a time, not trustworthiness when multiple pressures co-occur.

\label{sec:gap}These limitations motivate what we call the \emph{representativeness gap}. Let $\mathcal{S}$ denote the space of socio-technical scenarios relevant to agent deployment, $P_{\mathrm{deploy}}$ the real deployment distribution over $\mathcal{S}$, and $\mathbf{m}(a,s)=(m_{P_1}(a,s),\dots,m_{P_5}(a,s))$ a vector of per-property trustworthiness measurements for agent $a$ in scenario $s$. Deployment-level trustworthiness is then the property-wise expectation $\mathbf{T}(a) = \mathbb{E}_{s \sim P_{\mathrm{deploy}}}[\mathbf{m}(a,s)]$. Existing benchmarks estimate performance using samples from a benchmark distribution $P_{\mathrm{bench}}$ whose support is narrower and whose long-tail coverage is limited. When $P_{\mathrm{bench}}$ is misaligned with $P_{\mathrm{deploy}}$, benchmark scores become unreliable estimators of real-world trustworthiness, especially for low-frequency but high-consequence scenarios~\cite{eriksson2025trustbenchmarks}. From this perspective, the bottleneck of agent trustworthiness evaluation is no longer metric design alone, but \emph{distribution design}~\cite{liang2023helm,kiela2021dynabench,shirali2022dynamicbench,koh2021wilds}. Prior work expands metrics, tasks, or attack coverage; our claim is that the missing object is the \emph{representative, risk-sensitive scenario distribution itself}, jointly evaluated across all five properties.

\label{para:coverage_audit}
To make this gap concrete, we conduct an \emph{illustrative coverage audit} over a selected subset of prominent benchmarks (Table~\ref{tab:coverage}). We map each benchmark to eight evaluation dimensions and assign one of three coverage levels based on qualitative expert assessment: $\checkmark$ (primary: the dimension is a core evaluation target), $\triangle$ (partial: the dimension is indirectly or occasionally tested), or $\times$ (minimal: no meaningful coverage). The audit illustrates that current benchmarks are strongly \emph{axis-specialised}: they form valuable \emph{benchmark islands}, but collectively under-cover deployment-critical dimensions---especially operational constraints, social context, and consequence-sensitive risk---which correspond to properties P3--P5 in our definition.

\begin{table}[!htbp]
\centering
\caption{Coverage audit of selected benchmarks. $\checkmark$ = primary, $\triangle$ = partial, $\times$ = minimal (based on qualitative expert assessment).}
\label{tab:coverage}
\small
\setlength{\tabcolsep}{3pt}
\resizebox{\columnwidth}{!}{%
\begin{tabular}{@{}l cccc cccc@{}}
\toprule
& \multicolumn{4}{c}{\textit{Benchmark Slices}} & \multicolumn{4}{c}{\textit{Deployment Context}} \\
\cmidrule(lr){2-5} \cmidrule(lr){6-9}
\textbf{Benchmark} & \rotatebox{55}{Task} & \rotatebox{55}{Tool} & \rotatebox{55}{Long-Horiz.} & \rotatebox{55}{Factuality} & \rotatebox{55}{Adversarial} & \rotatebox{55}{Operational} & \rotatebox{55}{Social} & \rotatebox{55}{Risk} \\
\midrule
AgentBench      & $\checkmark$ & $\triangle$  & $\triangle$  & $\times$     & $\times$     & $\times$     & $\times$     & $\times$ \\
WebArena        & $\checkmark$ & $\checkmark$ & $\checkmark$ & $\times$     & $\times$     & $\times$     & $\times$     & $\times$ \\
SWE-bench       & $\checkmark$ & $\triangle$  & $\checkmark$ & $\times$     & $\times$     & $\triangle$  & $\times$     & $\times$ \\
HaluEval        & $\times$     & $\times$     & $\times$     & $\checkmark$ & $\times$     & $\times$     & $\times$     & $\times$ \\
JailbreakBench  & $\times$     & $\times$     & $\times$     & $\times$     & $\checkmark$ & $\times$     & $\times$     & $\triangle$ \\
API-Bank        & $\triangle$  & $\checkmark$ & $\triangle$  & $\times$     & $\times$     & $\times$     & $\times$     & $\times$ \\
AgentDojo       & $\triangle$  & $\checkmark$ & $\triangle$  & $\times$     & $\checkmark$ & $\triangle$  & $\times$     & $\triangle$ \\
\bottomrule
\end{tabular}}
\end{table}

\subsection{Provenance and Delta}
\label{subsec:provenance_delta}

The coverage audit in Table~\ref{tab:coverage} captures the broad benchmark landscape, but the closest adjacent works to HAAF are two tool-use evaluation suites: AgentDojo~\cite{debenedetti2024agentdojo} and ToolSandbox~\cite{lu2024toolsandbox}. To make the contribution of HAAF precise, we contrast all three along six axes that matter for deployment-level trustworthiness assessment (Table~\ref{tab:related_delta}).

\begin{table}[!htbp]
\centering
\caption{Adjacent tool-use benchmarks vs.\ HAAF. ``\# Properties'' counts deployment-level trustworthiness properties (P1--P5) explicitly scored; ``Red/Blue Cycle'' indicates whether the suite supports an iterative red-team $\to$ blue-team $\to$ re-evaluation loop; ``IR Scoring'' refers to Improper Refusal as a first-class evaluated class.}
\label{tab:related_delta}
\small
\setlength{\tabcolsep}{4pt}
\begin{tabular}{@{}l cccccc@{}}
\toprule
\textbf{Benchmark} & \textbf{\# Properties} & \textbf{\# Models} & \textbf{Red/Blue Cycle} & \textbf{IR Scoring} & \textbf{Sandbox} & \textbf{Scenario Source} \\
\midrule
AgentDojo~\cite{debenedetti2024agentdojo}     & 1 (injection only)        & 10 & no  & no  & 74 synthetic tools, 4 domains  & 97 tasks + 629 injections, hand-authored \\
ToolSandbox~\cite{lu2024toolsandbox}          & 0 (5 capability cats.)    & 13 & no  & no  & 34 synthetic tools, 11 domains & 1{,}032 cases, hand-crafted, seeded branching \\
HAAF (ours)                                   & 5 (P1--P5, incl.\ IR)     & 13 (7 families) & yes & yes & 5-tool synthetic sandbox       & 100 hand-authored scenarios (Layer-4 subset-selected) \\
\bottomrule
\end{tabular}
\end{table}

\paragraph{Provenance and delta.}
HAAF inherits the synthetic-sandbox tool-use evaluation paradigm pioneered by AgentDojo~\cite{debenedetti2024agentdojo} and ToolSandbox~\cite{lu2024toolsandbox}, but extends it along three axes that the adjacent benchmarks deliberately leave out. First, \emph{trustworthiness coverage}: AgentDojo scores a single property (indirect prompt-injection robustness) and ToolSandbox scores zero safety properties (its five categories are tool-use \emph{capability} dimensions), whereas HAAF scores all five deployment-level properties P1--P5 jointly on the same trajectory population. Second, \emph{model breadth}: both prior suites evaluate a fixed model list (10 and 13 respectively) at a single time slice; HAAF spans 13 models across 7 families with version-pair contrasts that expose anti-scaling regressions (\S\ref{subsec:antiscaling_stats}). Third, \emph{iterative protocol}: AgentDojo ships four static defenses and ToolSandbox has no defense loop, while HAAF closes a red-team $\to$ blue-team $\to$ re-evaluation cycle with paired-bootstrap significance testing. Finally, HAAF promotes Improper Refusal (IR) to a first-class evaluated class under P1; neither adjacent benchmark scores IR.
